\documentclass{article}
\usepackage[zh]{homework}

\config{分析\,-\,0}{2026 春季}{10}

\begin{document}

\begin{problem}
设 $f \in \mathcal{R}[a,b]$, 证明 $f$ 至少有一个连续点. (不能直接使用 Lebesgue 准则.)
\end{problem}

\begin{proof}
  反证法. 按照下一道题目的题干中的方法定义振幅 \( \omega_f(x) \). 则 \( \omega_f(x) \) 对于所有的 \( x \in [a,b] \) 都大于零. 设 \( \varepsilon=\min\{\omega_f(x):x\in[a,b]\}>0 \). 则对于任意的划分 \( P \), 均有
  \[ U(f,P)-L(f,P) = \sum_{i=1}^{n}\Delta x_i \omega_f(x_i) \ge \varepsilon (b-a)>0. \]

  这与 \( f \in \mathcal{R}[a,b] \) 矛盾.
\end{proof}



\begin{problem}
设 $f:[a,b] \to \R$ 有界. 定义 $f$ 在 $x \in [a,b]$ 的振幅为
\[
\omega_f(x) = \lim_{\delta \to 0^+} \sup\{\abs{f(s) - f(t)} : s,t \in N_\delta(x) \cap [a,b]\}.
\]
给定 $\varepsilon > 0$, 定义
\[
D_\varepsilon = \{x \in [a,b] : \omega_f(x) \ge \varepsilon\}.
\]
\begin{enumerate}[label=(\alph*)]
	\item 证明 $\omega_f(x)$ 定义中的极限一定存在.
	\item 证明 $f$ 在 $x$ 处连续当且仅当 $\omega_f(x)=0$.
	\item 证明 $D_\varepsilon$ 是闭集.
	\item 设 $f \in \mathcal{R}[a,b]$, 证明 $D_\varepsilon$ 是零测集.
	\item 令 $D_f$ 是 $f$ 在 $[a,b]$ 上的间断点集, 设 $D_f$ 是零测集, 证明 $f \in \mathcal{R}[a,b]$.
\end{enumerate}
\end{problem}

\begin{proof}
  (a) 注意到, 对于 \( 0<\delta<\delta' \), 有 \( (N_\delta(x)\cap[a,b])\subset  (N_{\delta'}(x)\cap[a,b]) \), 因而 \( \sup\{\abs{f(s) - f(t)} : s,t \in N_\delta(x) \cap [a,b]\} \) 关于 $\delta$ 单调递减, 且有下界 \( 0 \), 故极限存在.

  (b) 若 \( f \) 在 \( x \) 处连续, 则根据定义, 对于任意的 \( \varepsilon>0 \), 存在 \( \delta>0 \), 使得对于任意的 \( t\in N_\delta(x)\cap[a,b] \), 都有 \( \abs{f(t)-f(x)}<\varepsilon/2 \). 从而对于任意的 \( s,t\in N_\delta(x)\cap[a,b] \), 都有 \( \abs{f(s)-f(t)}<\varepsilon \). 因此\( \omega_f(x)=0 \). 

  反之, 若 \( \omega_f(x)=0 \), 则对于任意的 \( \varepsilon>0 \), 存在 \( \delta>0 \), 使得对于任意的 \( s,t\in N_\delta(x)\cap[a,b] \), 都有 \( \abs{f(s)-f(t)}<\varepsilon \). 取 \( s=x \), 则对于任意的 \( t\in N_\delta(x)\cap[a,b] \), 都有 \( \abs{f(t)-f(x)}<\varepsilon \). 因此\( f \) 在 \( x \) 处连续.

  (c) 只需证明 \( D_\varepsilon \) 的补集
  \[ D_\varepsilon^c = \{x \in [a,b] : \omega_f(x) < \varepsilon\} \]
  是开集. 对于 \( x_0\in D_\varepsilon^c \), 由于 \( \omega_f(x_0)<\varepsilon \), 则存在 \( \delta_0 \), 使得
  \[ \sup\{\abs{f(s) - f(t)} : s,t \in N_{\delta_0}(x_0) \cap [a,b]\}=:\varepsilon_0 < \varepsilon. \]

  我们来证明: 对于任意的 \( x\in N_{\delta_0}(x_0)\cap[a,b] \), 都有 \( \omega_f(x)\le \varepsilon_0<\varepsilon \). 取 \( \delta=\delta_0-\abs{x-x_0}>0 \) (因为 \( x\in N_{\delta_0}(x_0) \)), 则不难知道 \( (N_\delta(x)\cap[a,b])\subset (N_{\delta_0}(x_0)\cap[a,b]) \). 从而
  \[ \sup\{\abs{f(s) - f(t)} : s,t \in N_\delta(x) \cap [a,b]\} \le \sup\{\abs{f(s) - f(t)} : s,t \in N_{\delta_0}(x_0) \cap [a,b]\}=\varepsilon_0. \]
  因此 \( \omega_f(x)\le \varepsilon_0<\varepsilon \), 从而 \( x\in D_\varepsilon^c \). 这说明 \( N_{\delta_0}(x_0)\cap[a,b]\subset D_\varepsilon^c \), 从而 \( D_\varepsilon^c \) 是开集, \( D_\varepsilon \) 是闭集.

  (d) 对于任意的正实数 \( \eta \), 由于 \( f \in \mathcal{R}[a,b] \), 则存在划分 \( P \), 使得 
  \[ \sum_{i=1}^{n}\Delta x_i \omega_i < \frac{1}{2}\varepsilon\eta ,\]
  其中
  \[ \omega_i = \sup\{\abs{f(s) - f(t)} : s,t \in [x_{i-1}, x_i]\}. \]

  对于 \( 1\le i\le n \), 记闭区间 \( E_i=[x_{i-1},x_i] \), 开区间 \( O_i=\left( x_{i-1}-\frac{1}{2}\Delta x_i, x_{i}+\frac{1}{2}\Delta x_i \right) \). 则 \( E_i\subset O_i \), 且 \( |O_i|=2 \Delta x_i \). 记
  \[ \left\{ i: \exists\, x\in E_i, \omega_f(x)\ge \varepsilon \right\}=:I\subset J:=\left\{ j: \omega_j\ge \varepsilon \right\}. \]

  则 \( D_\varepsilon \subset \bigcup_{i\in I} E_i \subset \bigcup_{i\in I} O_i\subset \bigcup_{j\in J} O_j \). 由于
  \[ \frac{1}{2}\eta \varepsilon>\sum_{i=1}^{n}\Delta x_i \omega_i\ge \sum_{j\in J} \frac{1}{2}|O_j|\varepsilon
  \implies \sum_{j\in J}^{}|O_j| < \eta, \]
  再结合 \( \eta \) 的任意性, 知 \( D_\varepsilon \) 可以被总长度任意小的开区间覆盖. 因此 \( D_\varepsilon \) 是零测集.

  (e) 由于任意的 \( \varepsilon>0 \), 都有 \( D_\varepsilon\subset D_f \), 则 \( D_\varepsilon \) 是零测集. 任取 \( \eta>0 \), 则对于任意的 \( \varepsilon>0 \), 都存在可数的开区间集合 \( \{O_i\} \), 使得 \( D_\varepsilon\subset \bigcup_i O_i \) 且 \( \sum_i |O_i|<\eta \). 而由 \( D_\varepsilon \) 是闭的, 因此存在有限的开区间集合 \( \{O_i\}_{i=1}^n \), 使得 \( D_\varepsilon\subset \bigcup_{i=1}^n O_i \) 且 \( \sum_{i=1}^n |O_i|<\eta \). 

  由于 \( [a,b]- D_\varepsilon \) 上任意一点关于 \( f \) 的振幅都小于 \( \varepsilon \), 因此对于任意的 \( x\in [a,b]-\bigcup_{i=1}^n O_i \), 存在 \( x \) 的开邻域 \( U(x) \), 使得对于任意的 \( s,t\in U(x)\cap[a,b] \), 都有 \( \abs{f(s)-f(t)}<\varepsilon \). 记 \( E= [a,b]-\bigcup_{i=1}^n O_i \), 则 \( E \) 为闭集, 并且 \( E\subset\bigcup_{x\in E}U(x) \), 从而存在子覆盖, 设为 \( \bigcup_{j=1}^{m}U_j \). 

  考虑所有 \( O_i \) 与 \( U_j \) 的端点, 将这些点视为划分点, 构成一个划分 \( P:a=x_0<x_1<\dots<x_N=b \). 设 \( M=\sup\{\abs{f(s)-f(t)}:s,t\in[a,b]\} \). 则对于任意的 \( 1\le k\le N \), 均有 \( \omega_k\le M \). 
  
  从而
 \[ U(f,P)-L(f,P) = \sum_{j}\Delta x_j \omega_j 
 \le \abs{\bigcup_{i=1}^{n}O_i}M + \left( (b-a)-\abs{\bigcup_i O_i} \right)\varepsilon
 <\eta M+(b-a)\varepsilon. \]

 由 \( \varepsilon \) 与 \( \eta \) 的任意性, 知 \( \eta M+(b-a)\varepsilon \) 可以任意小. 因此 \( f \in \mathcal{R}[a,b] \).
\end{proof}



\begin{problem}
设 $f$ 在 $[a,b]$ 上满足 $f \ge 0$, $f$ 连续且 $\int_a^b f(x)\,\d x = 0$. 证明对一切 $x \in [a,b]$, $f(x)=0$.
\end{problem}

\begin{proof}
  采用反证法. 否则, 设 \( x_0\in [a,b] \) 使得 \( f(x_0)>0 \). 由于 \( f \) 连续, 存在 \( x_0 \) 的开邻域 \( N_\delta(x_0) \), 使得对于任意的 \( x\in N_\delta(x_0)\cap[a,b] \), 都有 \( \abs{f(x)-f(x_0)}<\frac{1}{2}f(x_0) \), 即 \( f(x)>\frac{1}{2}f(x_0) \). 从而
  \[ \int_a^b f(x)\,\d x \ge \int_{N_\delta(x_0)\cap[a,b]} f(x)\,\d x > \frac{1}{2}f(x_0)|N_\delta(x_0)\cap[a,b]|>0, \]
  矛盾.
\end{proof}



\begin{problem}
设 $f:[a,b] \to \R$ 有界.
\begin{enumerate}[label=(\alph*)]
	\item 若已知 $f^2 \in \mathcal{R}[a,b]$, 是否必然 $f \in \mathcal{R}[a,b]$? 证明或举出反例.
	\item 若已知 $f^3 \in \mathcal{R}[a,b]$, 答案是否改变?
\end{enumerate}
\end{problem}

\begin{solution}
  (a) 不必然. 反例: 取
  \[ f(x)=\begin{cases}
  1, & x\in \Q\cap[a,b], \\
  -1, & x\in ([a,b]\cap\R)-\Q.
  \end{cases} \]
  则 \( f^2(x)=1 \) 对所有 \( x\in [a,b] \) 成立, 因此 \( f^2 \in \mathcal{R}[a,b] \). 但 \( f \) 在 \( [a,b] \) 上处处不连续, 因此 \( f \notin \mathcal{R}[a,b] \).

  (b) 改变. 这是因为映射 \( x\mapsto \sqrt[3]{x} \) 在 \( \R \) 上连续, 所以如果 \( f^3 \in \mathcal{R}[a,b] \), 则 \( f = \sqrt[3]{f^3} \in \mathcal{R}[a,b] \).
\end{solution}



\begin{problem}
设 $f$ 和 $g$ 在 $[a,b]$ 上均有界, $f \in \mathcal{R}[a,b]$.
\begin{enumerate}[label=(\alph*)]
	\item 若 $f$ 和 $g$ 只在有限个点上取值不同, 证明 $g \in \mathcal{R}[a,b]$, 且 $\int_a^b f(x)\,\d x = \int_a^b g(x)\,\d x$.
	\item 若 $f$ 和 $g$ 在可数无穷多个点上取值不同, 结论是否成立? 证明或举出反例.
\end{enumerate}
\end{problem}

\begin{solution}
  (a) 设 \( h(x)=g(x)-f(x) \). 则由题目条件, \( h \) 仅在有限个点 \( x_1, \dots,x_n \) 上取值非零. 设 \( M=\max\{\abs{h(x_i)}:1\le i\le n\} \). 对于充分大的正整数 \( N \), 取 \( [a,b] \) 的等距划分 \( P: a=x_0<x_1<\dots<x_N=b \). 则至多只有 \( n \) 个子区间 \( [x_{i-1}, x_i] \) 上 \( h \) 取值不恒为 \( 0 \), 从而
  \[ \abs{\sum_{i=1}^{N}f(\xi_i)\Delta x_i}\le \frac{b-a}{N}\sum_{i=1}^{n}f(x_i)=\frac{(b-a)M}{nN}. \]

  由于 \( N \) 可以充分大, 从而 \( h \in \mathcal{R}[a,b] \), 并且 \( \int_{a}^{b}h(x)\d x=0 \). 因此 \( g=f+h \in \mathcal{R}[a,b] \), 且 
  \[ \int_a^b f(x)\,\d x = \int_a^b g(x)\,\d x. \]

  (b) 不成立. 反例: 取 \( f \) 恒为 \( 0 \), 而 \( g \) 在 \( \Q\cap[a,b] \) 上取值 \( 1 \), 在 \( ([a,b]\cap\R)-\Q \) 上取值 \( 0 \). 则 \( f \in \mathcal{R}[a,b] \), 且 \( g \) 与 \( f \) 在可数无穷多个点上取值不同. 但 \( g \) 在 \( [a,b] \) 上处处不连续, 因此 \( g \notin \mathcal{R}[a,b] \).
\end{solution}



\begin{problem}
计算下列定积分或不定积分:
\begin{enumerate}[label=(\alph*)]
	\item $\int_1^e \frac{1}{x}\,\d x$.
	\item $\int (3x^2 - 2x + 1)\,\d x$.
	\item $\int_{-1}^2 \abs{x - 1}\,\d x$.
	\item 设
	\[
	f(x) =
	\begin{cases}
		x^2, & 0 \le x < 1, \\
		2 - x, & 1 \le x \le 2.
	\end{cases}
	\]
	求 $\int_0^2 f(x)\,\d x$.
\end{enumerate}
\end{problem}

\begin{solution}
  (a) \[ \int_1^e \frac{1}{x}\,\d x = \ln x \Big|_1^e = 1. \]

  (b) \[ \int (3x^2 - 2x + 1)\,\d x = x^3 - x^2 + x + C. \]

  (c) \[ \int_{-1}^2 \abs{x - 1}\,\d x = \int_{-1}^1 (1-x)\,\d x + \int_1^2 (x-1)\,\d x = \frac{5}{2}. \]

  (d) \[ \int_0^2 f(x)\,\d x = \int_0^1 x^2\,\d x + \int_1^2 (2-x)\,\d x = \frac{1}{3} + \frac{1}{2} = \frac{5}{6}. \]
\end{solution}



\begin{problem}
设 $f \in \mathcal{R}[a,b]$, 令 $A = \frac{1}{b-a}\int_a^b f(x)\,\d x$, 证明存在 $\xi, \eta \in [a,b]$ 使得
\[
f(\xi) \ge A, \quad f(\eta) \le A.
\]
提示: 可以利用第一题的结论.
\end{problem}

\begin{proof}
  由 \( A \) 的定义, 有 \( \int_{a}^{b}(f(x)-A)\d x=0 \). 若不存在 \( \xi\in[a,b] \) 使得 \( f(\xi)\ge A \), 则 \( f-A \) 在 \( [a,b] \) 上恒负, 因而在 \( [a,b] \) 上的积分为 \( 0 \), 矛盾. 同理, 若不存在 \( \eta\in[a,b] \) 使得 \( f(\eta)\le A \), 则 \( f-A \) 在 \( [a,b] \) 上恒正, 因而在 \( [a,b] \) 上的积分为 \( 0 \), 矛盾.

因此存在 \( \xi, \eta\in[a,b] \) 使得 \( f(\xi)\ge A \) 且 \( f(\eta)\le A \).
\end{proof}



\begin{problem}
考虑
\[
I(\lambda) = \int_a^b f(x)\cos \lambda x\,\d x.
\]
\begin{enumerate}[label=(\alph*)]
  \item 设 $f(x)$ 在 $[a,b]$ 上单调有界, 证明 $\lim_{\lambda \to \infty} I(\lambda) = 0$.
  \item 设 $f \in C([a,b])$, 证明 $\lim_{\lambda \to \infty} I(\lambda) = 0$.
\end{enumerate}
\end{problem}

\begin{proof}
  我们直接来证明 (b). 
\end{proof}



\begin{problem}
计算下列定积分或不定积分.
\begin{enumerate}[label=(\alph*)]
  \item $\int \frac{x+1}{x^2 - 5x + 6}\,\d x$.
  \item $\int \frac{x^3 + x^2 + x + 1}{x^2 + 1}\,\d x$.
  \item $\int x^2 \ln x\,\d x$.
  \item $\int \sin^3 x \cos^2 x\,\d x$.
  \item $\int \cos^4 x\,\d x$.
  \item $\int_0^{\pi/2} \frac{1}{\cos x + \sin x}\,\d x$.
\end{enumerate}
\end{problem}

\begin{problem}
假设 $f \in C^1([a,b])$, $f(a)=f(b)=0$, 并且
\[
\int_a^b f^2(x)\,\d x = 1.
\]
证明
\[
\int_a^b x f(x)f'(x)\,\d x = -\frac{1}{2}
\]
和
\[
\int_a^b [f'(x)]^2\,\d x \cdot \int_a^b x^2 f^2(x)\,\d x > \frac{1}{4}.
\]
提示: 利用第 9 次作业中证明的 Schwarz 不等式.
\end{problem}

\begin{problem}
设 $f \in \mathcal{R}[a,b]$. 定义
\[
\lVert f \rVert_1 = \int_a^b \abs{f(x)}\,\d x, \quad
\lVert f \rVert_2 = \left\{\int_a^b \abs{f(x)}^2\,\d x\right\}^{1/2}.
\]
给定 $\varepsilon > 0$.
\begin{enumerate}[label=(\alph*)]
  \item 证明在 $[a,b]$ 上存在连续函数 $g$ 满足 $\lVert f-g \rVert_2 < \varepsilon$.
  \item 证明在 $[a,b]$ 上存在连续函数 $h$ 满足 $\lVert f-h \rVert_1 < \varepsilon$.
\end{enumerate}
提示: 设 $P = \{x_0,\dots,x_n\}$ 是 $[a,b]$ 的一个适当的分割, 若 $x_{i-1} < t < x_i$, 定义
\[
g(t) = \frac{x_i - t}{\Delta x_i} f(x_{i-1}) + \frac{t - x_{i-1}}{\Delta x_i} f(x_i).
\]
\end{problem}

\begin{problem}
定义
\[
f(x) = \int_x^{x+1} \sin(t^2)\,\d t.
\]
\begin{enumerate}[label=(\alph*)]
  \item 求证 $x > 0$ 时 $\abs{f(x)} < \frac{1}{x}$.
  提示: 置 $t^2 = u$, 再分部积分以证 $f(x)$ 等于
  \[
  \frac{\cos(x^2)}{2x} - \frac{\cos[(x+1)^2]}{2(x+1)} - \int_{x^2}^{(x+1)^2} \frac{\cos u}{4u^{3/2}}\,\d u,
  \]
  用 $-1$ 代替 $\cos u$.
  \item 证明
  \[
  2x f(x) = \cos(x^2) - \cos[(x+1)^2] + r(x),
  \]
  其中 $\abs{r(x)} < c/x$, 而 $c$ 是常数.
  \item 求 $x f(x)$ 当 $x \to \infty$ 时的上, 下极限.
\end{enumerate}
\end{problem}



\end{document}